\section{Cross-Pathway Synthesis}
\label{sec:synthesis}

Sections~\ref{sec:tok2tok}--\ref{sec:composite} have surveyed, pathway by pathway, the mechanisms the community uses to re-represent agent experience from one carrier into another. But the problems faced in building an agent system cut naturally across these pathways: given one trajectory, should it go through Narrative Abstraction (P1) into a reflection for retrieval, or through Policy Internalization (P5) into the weights for zero marginal inference cost? Once the decision is made, the experience's interpretability, editability, inference efficiency, and maintenance cost are locked into a particular trade-off profile, and later correction is expensive. The existing literature, divided by component (memory, planning, tool use) or by technique (RAG, SFT, RLHF), offers no analytical tool for this cross-pathway choice---each gives depth within a single carrier or a single pathway, but no framework tells you under what conditions one pathway is better than another, and why.

This section steps back from the method-level detail of Sections~\ref{sec:tok2tok}--\ref{sec:composite} and synthesizes structurally at the pathway level and carrier level. It introduces no new literature but uses the existing evidence to establish three generative principles (Section~\ref{sec:principles}), uses them to explain in a unified way the differences in production cost across the seven pathways (Section~\ref{sec:pathway-tradeoff}) and the differences in consumption characteristics across the five carriers (Section~\ref{sec:utilization-analysis}), and finally crosses these differences under the lens of task constraints to derive the logic of utilization-driven pathway selection (Section~\ref{sec:pathway-selection}).

\subsection{Three Generative Principles}
\label{sec:principles}

Section~\ref{sec:carriers} defined the position of the five carriers on the Tokenized $\rightarrow$ Latent $\rightarrow$ Parametric spectrum, and Section~\ref{sec:transformation} defined the source and target ends of the seven pathways. But carrier classification and pathway definition answer ``what,'' not ``why''---why is the consumption cost of Narrative a context token paid per use, while Policy is zero-marginal? Why do the two Tokenized $\rightarrow$ Tokenized pathways need only LLM inference, while the three Tokenized $\rightarrow$ Parametric pathways all need gradient-based training? Why do composite pathways (Section~\ref{sec:composite}) form mainly around Policy, while Evaluator is almost always a terminus?

These are not isolated empirical observations. Behind them are three structural constraints---the carrier's form of existence, its degree of formalization, and its functional role---each locking the behavior of that carrier and its production pathway along some dimension. This section makes the three constraints explicit as generative principles, the unified coordinate system for the analysis that follows.

\subsubsection{Existence-Level Principle}
\label{sec:existence-principle}

\begin{quote}\itshape
A carrier's position on the Tokenized $\rightarrow$ Latent $\rightarrow$ Parametric spectrum simultaneously determines its consumption interface, per-use cost, control freedom, and traceability. These dimensions do not vary independently---they are locked monotonically by the existence level.
\end{quote}

The derivation starts from the agent decision rule $a_t \sim \pi_\theta(\cdot \mid c_t)$. There are only two structural possibilities for how experience is consumed: it enters $c_t$ (external, explicit, participating in the forward pass of the current context), or it enters $\theta$ (internal, implicit, fixed in the model parameters or intermediate states). The carrier's existence level decides which path it takes.

A Tokenized carrier exists as a discrete symbol outside the model. Its consumption interface is necessarily context injection---the experience must first be retrieved, then encoded into $c_t$. Its per-use cost therefore includes retrieval latency and the prefill overhead of context processing, and it worsens as the volume of experience grows (the context window has limited capacity, and retrieval precision falls as the store grows). The other side of a discrete symbol is independent addressability: each reflection, each rule can be located, read, edited, or deleted on its own, without touching other entries. Control freedom is high, and traceability is high---you can ask which experience a given decision drew on and get a readable answer.

A Parametric carrier exists as continuous weights inside the model. Its consumption interface is necessarily forward-pass computation---the weights are always present, the forward pass activates them automatically, with no retrieval and no extra context token. The per-use cost is zero-marginal. But the price of weights is the loss of independent addressing: a single experience's contribution is submerged in the accumulated batch gradient, and one cannot selectively activate ``this experience and not that one'' at inference. Control freedom is near zero---you cannot add, remove, or change any experience without retraining. Traceability is near zero---you cannot ask which training trajectory a decision used.

A Latent carrier sits between the two. Its consumption interface is attention integration within the forward pass---no extra token encoding, but a choice of which latent bank or which KV caches to load. Its per-use cost is low but not zero (the attention computation has its own cost, as does retrieving the latent bank). A continuous vector cannot be read directly, but it can be loaded selectively (which session, which bank). Control freedom and traceability are intermediate.

Along the Tokenized $\rightarrow$ Latent $\rightarrow$ Parametric spectrum, Figure~\ref{fig:spectrum} traces a single monotone curve.

\begin{figure}[t]
\centering
\fbox{%
\begin{minipage}{0.9\columnwidth}
\centering
\textbf{Tokenized} $\;\longrightarrow\;$ \textbf{Latent} $\;\longrightarrow\;$ \textbf{Parametric}\\[4pt]
\emph{Interpretability, editability, traceability, control}: decreasing $\longrightarrow$\\[2pt]
\emph{Inference efficiency, capacity scalability}: increasing $\longrightarrow$
\end{minipage}}
\caption{The carrier spectrum. Along Tokenized $\rightarrow$ Latent $\rightarrow$ Parametric, interpretability, editability, traceability, and control fall monotonically while inference efficiency and capacity scalability rise; the two groups move in opposite directions and are locked together by existence level.}
\label{fig:spectrum}
\end{figure}

This is not an empirical summary of the five carriers; it is a structural corollary of the existence level. It yields a strong negative statement: \textbf{no carrier can hold both zero marginal inference cost and high interpretability at once.} A scheme that claims to be ``fast and transparent'' must be vague about its carrier's existence level---it is either in fact maintaining a Tokenized copy for audit (that is two carriers, not one), or its ``transparency'' refers to some approximate explanation (probing, attribution) rather than genuine readability.

The existence-level principle runs through all the analysis that follows. For Section~\ref{sec:pathway-tradeoff} (the production side), whether a transformation pathway crosses existence levels (P3/P4/P5/P7) or stays within one (P1/P2/P6) determines the scale of production cost and information loss more fundamentally than the choice of target carrier. For Section~\ref{sec:utilization-analysis} (the consumption side), the consumption profiles of the five carriers are not five independent data points but five sampling points along the same curve, indexed by spectrum position. For Section~\ref{sec:pathway-selection} (the selection side), utilization-driven selection is at bottom a selection of existence level---first decide which region of the spectrum the experience should sit in, then choose the specific carrier subclass within that region.

\subsubsection{Formalization Principle}
\label{sec:formalization-principle}

\begin{quote}\itshape
Within the Tokenized level, the difference in formalization between Narrative and Schematic decides whether experience is consumed by understanding or consumed mechanically, and in turn places the verifiability of the product in different categories---not different degrees of ``high verifiability.''
\end{quote}

The existence-level principle explains the difference between Tokenized and Parametric. But the Tokenized level has two subclasses---Narrative ($\mathcal{C}^T_N$) and Schematic ($\mathcal{C}^T_S$)---that share the context-injection consumption interface yet differ in verifiability in a same-but-different way. This difference cannot be explained by the existence-level principle (both are at the same level) and needs a second principle.

Narrative is consumed through language and multimodal understanding. An agent injects a reflection, a rule, or a summary into the context and relies on the model's language understanding to extract decision guidance from it. Its verification can only be by reading---a human reviewer or a strong LLM judges whether the content is correct, relevant to the current task, and consistent with existing experience. This verification is semantic; it covers ``whether it is said right.'' But it has three inherent limits: (i) it is subjective---two reviewers may judge the quality of the same reflection differently; (ii) it does not guarantee execution---a reflection judged ``correct'' may not be usable effectively by the agent in practice, ``said right but cannot be done''; (iii) it relies on language-understanding ability, and the agent's language understanding at inference may differ from the verifier's.

Schematic is consumed through parsing, execution, and traversal. A code skill is executed by an executor, a workflow DAG is traversed by an engine, a knowledge graph is retrieved by a graph query. Their verification is by running---an execution test gives an objective pass/fail signal, parsing gives an objective syntax-correct/error signal. This verification is formal; it covers ``whether it runs.'' It does not depend on subjective judgment, and test coverage is quantifiable and scalable. But its blind spot is exactly Narrative verification's strength: a code skill that ``runs'' may implement the wrong functionality, and a syntactically correct workflow may encode the wrong task decomposition.

The two verifications cover different categories of error. Narrative verification misses ``said well but cannot be done''; Schematic verification misses ``runs but points the wrong way.'' Both are ``high verifiability,'' but not the same verifiability---a difference of category, not of degree.

The analytical consequences of the formalization principle enter the pathway comparison of Section~\ref{sec:pathway-tradeoff} and the safety-constraint analysis of Section~\ref{sec:pathway-selection} directly. For Section~\ref{sec:pathway-tradeoff}: P1 and P2 are both Tokenized $\rightarrow$ Tokenized same-level transformations, yet their Output Verifiability is heterogeneous---P1's verification is the subjective judgment of ``reading,'' P2's the objective test of ``running.'' For Section~\ref{sec:pathway-selection}: when a task's safety-criticality comes from ``must not do the wrong thing'' (execution safety), one should lean toward Schematic---an error can be blocked before deployment through a testing gate; when safety-criticality comes from ``must not understand wrong'' (semantic safety), one should lean toward Narrative---a human can review the content before the experience enters the context.

\subsubsection{Functional-Role Principle}
\label{sec:functional-role-principle}

\begin{quote}\itshape
Within the Parametric level, the difference in functional role between Policy and Evaluator decides the carrier's topological position in the pathway graph---Policy is a hub node, Evaluator a sink node. This is not an empirical observation but a graph-theoretic property entailed directly by the definition of the functional roles.
\end{quote}

The existence-level principle explains the difference between the Parametric level and the Tokenized level. But the Parametric level has two subclasses of sharply different functional role---Policy ($\mathcal{C}^P_\pi$) and Evaluator ($\mathcal{C}^P_\phi$). Both have forward-pass computation as their consumption interface, zero-marginal cost, and near-zero traceability. Yet they diverge on one key dimension: where they can travel in the pathway graph.

Policy's functional role is to produce actions. An action, executed in the environment, produces a new trajectory (a sequence of observations). A trajectory is a Tokenized carrier---which means every use of the Policy automatically produces a product consumable by any pathway sourced in Tokenized (P1/P2/P4/P5/P7). The Policy externalizes explicitly through P7 (Parametric Externalization) and implicitly through rollout, re-entering the whole pathway graph. In the directed graph of the seven pathways, the Policy can exit through P7 and return through P5, forming a loop---it is a member of a strongly connected component.

Evaluator's functional role is to produce evaluation signals---a scalar score, a preference label, a critique text. An evaluation signal has only one meaningful mode of consumption: to update, select, or guide a Policy. Without combining with a Policy (P6), the evaluation signal has nowhere to go. In the pathway graph, the Evaluator has only the in-edge P4 (Evaluator Internalization) and the out-edge P6 (Evaluator-Driven Optimization). It is a sink node.

This principle is not empirical---not a matter of the community not yet exploring other downstream uses of the Evaluator, but that the very semantics of an evaluation signal is ``a judgment of behavior,'' and its natural consumer is the producer of behavior. One can imagine using an Evaluator signal to update another Evaluator (that is incremental training of P4, not a new pathway), or using an Evaluator signal to filter Tokenized experience (that is the Evaluator used as a filter in the preprocessing of P1/P5, not the establishment of an independent Evaluator $\rightarrow$ Tokenized transformation). Neither is the Evaluator entering a new pathway as a source; both are the Evaluator assisting an existing pathway as a tool. An independent new pathway requires the Evaluator's experience to be re-represented as another carrier---which is equivalent to turning an evaluation criterion into a generation capability, and the only known mechanism for that is P6.

The functional-role principle explains a salient pattern in Section~\ref{sec:composite}: composite pathways are built mainly around the Policy. Among the three composition patterns identified there, the integration of Evaluator--Policy Co-Evolution (Section~\ref{sec:coevolution}) acts on the coupling between Evaluator and Policy, Refinement-Mediated Policy Internalization (Section~\ref{sec:refinement-mediated}) ends at the Policy, and Generative Experience Curation (Section~\ref{sec:experience-curation}) centers on the Policy $\rightarrow$ Tokenized $\rightarrow$ Policy self-generation loop. This is not an accident of research taste---the Policy's hub position in the graph makes it the natural anchor node for composite pathways. The Evaluator is the opposite: once experience reaches the Evaluator, its journey in the pathway graph is over (the only exit is into the Policy). Choosing the Evaluator as the target carrier means accepting the terminalization of experience management.

The three principles cover three orthogonal axes of the carrier space: existence level (the position of experience in the architecture), degree of formalization (within the Tokenized level, how experience is consumed), and functional role (within the Parametric level, the experience's topological fate in the pathway graph). Each carrier's coordinates on these three axes---not its name---determine its production-cost profile (Section~\ref{sec:pathway-tradeoff}) and its consumption-cost profile (Section~\ref{sec:utilization-analysis}). Every analytical dimension in the next two sections traces back to at least one of these three principles.

\subsection{Pathway Trade-off Comparison}
\label{sec:pathway-tradeoff}

The object of pathway comparison is the structural property of the transformation operation itself, not the static attribute of the target carrier. This boundary is easy to cross in writing yet must be held: the goal here is not to compare ``which is more interpretable, Policy or Narrative''---that is the carrier-attribute comparison already done in Section~\ref{sec:carriers}. The goal is to compare the differences between the transformation from trajectory to reflection (P1) and the transformation from trajectory to policy weights (P5) in information loss, production cost, and incremental capability. The former compares carriers, the latter compares transformations---a different angle, a different conclusion.

For example, ``Policy has low interpretability'' is an attribute of the Policy carrier; used to describe P5's trade-off, it applies equally to P4 (also a Parametric product) but not to P6 (whose product is also Policy, but with an entirely different cost profile). The trade-off that genuinely belongs to P5 is that it needs large-scale trajectories as input, performs the transformation through gradient-based training, may lose information through underfitting or reward hacking rather than deliberate abstraction, produces a result untraceable to any single source experience, and requires retraining to add new experience. These are properties P1 lacks entirely---not a difference of degree but a structural difference.

Below, the analysis is organized into seven dimensions across the three cross-sections of the transformation operation: input side, process side, and output side. Dimensions within one cross-section share a viewpoint, which makes cross-pathway comparison easier.

\subsubsection{Input-Side Dimensions}

\paragraph{Source Carrier.} The source-carrier type comes straight from the pathway definition. The source ends of the seven pathways fall into two sets: Tokenized-sourced are P1/P2/P3/P4/P5 (five), and Parametric-sourced are P6/P7 (two). The common feature of a Tokenized source is that the input is readable, auditable, and filterable item by item---you can check the quality of each trajectory before deciding whether to use it. The common feature of a Parametric source is that the input is opaque---P6 sources from Evaluator weights, P7 from Policy weights, and neither's internal experience distribution can be inspected directly, so transformation quality depends heavily on the quality and calibration state of the source weights. This dimension is constant; it is listed not to compare but to establish an identity for each row.

\paragraph{Data Prerequisite.} The seven pathways' requirements on input data differ across three sub-dimensions.

\textbf{Volume.} The granularity of a single trajectory separates lightweight transformation from retraining transformation. P1 (single-trajectory variants such as Reflexion~\cite{Shi23b}) and P2 (single-skill generation such as Voyager~\cite{Wan23c}) can run immediately at the end of one trial, without waiting for experience to accumulate. P1's cross-trajectory variants (such as AutoGuide~\cite{Fu24} and ExpeL~\cite{Zha23c}) and P7 need multiple trajectories. P3 (Cache-Based) can be built from a single session. P3 (Prompt/Module-Based), P4, and P5 need large-scale training data---from hundreds to millions of trajectories. P6's data-volume need depends on the type and deployment of the Evaluator: online RL (such as OAIF~\cite{Guo24}) updates continuously on the rollout stream, the trajectory volume set by the number of training steps; offline DPO needs a pre-built set of preference pairs.

\textbf{Quality Signal.} The type of supervision signal decides which experiences can be transformed. The basic gradient is as follows:
\begin{itemize}
  \item \emph{Trajectory alone} (no explicit supervision): P1 and P2 can extract an insight, skill, or workflow from unlabeled trajectories, requiring no success/failure label. A failed trajectory can expose an error pattern through post-hoc analysis (such as Reflexion's failure reflection), and a successful trajectory can induce a reusable strategy.
  \item \emph{Outcome-level signal} (weak supervision): P4's outcome-supervised Evaluator needs a trajectory-level success/failure judgment or pairwise preference. P5 (SFT variants) usually needs filtered successful trajectories or high-quality demonstrations. P7's generation quality often relies on outcome filtering at sampling time.
  \item \emph{Process-level signal} (process supervision): P4's process-supervised Evaluator (such as a PRM) needs step-level correctness annotation. P6's discriminative process-evaluator variants need step-level value estimates or quality scores.
  \item \emph{Human preference annotation} (strong supervision): the standard paradigm of P4's RM training needs human preference labels. P6's canonical RLHF pipeline shares this need.
  \item \emph{Success signal plus environment execution feedback}: P2's Programmatic Skill Construction needs, beyond trajectories, an execution-verification loop---a code skill is admitted only after it actually runs in the environment and passes the tests.
\end{itemize}

\textbf{Source Diversity.} The experience source limits the scope of a transformation's applicability. P1/P2/P3/P5/P7 can use agent self-generated trajectories as well as human demonstrations or teacher-model-synthesized data. P4's RM training usually relies on human labels, though AI-generated preferences (such as Constitutional AI~\cite{Bai22b}) reduce the dependence on human annotation step by step. P6's Evaluator-to-Policy stage in principle does not touch the original experience source---its input is the signal the Evaluator produces, the diversity of the experience source already folded into the Evaluator weights back in the P4 stage.

\paragraph{Input Robustness.} Transformation operations differ widely in their tolerance for noisy input. P1's design naturally includes noise handling: the reflection mechanism extracts diagnostic information from failed trajectories specifically, and the cross-trajectory induction mechanism filters out incidental factors through multi-trajectory comparison---noise is exactly P1's input material. P2's Programmatic Skill Construction filters noise through execution verification: an unqualified skill is eliminated at the test stage and never admitted. P3 (Cache-Based) saves the KV produced by the forward pass directly and does not actively denoise---the KV of a noisy trajectory is saved just the same and may be recalled at retrieval.

P5 (SFT) is sensitive to noise. A low-quality trajectory in the training data directly lowers policy quality. The community's practice is to pre-filter---using only successful or high-scoring trajectories for SFT---but the filtering criterion itself may introduce bias (keeping only trajectories of a certain style). P5 (RL) amplifies the sensitivity to noise further: reward hacking is precisely the policy exploiting an unintended noise pattern in the reward signal. P4's annotation noise (annotator inconsistency, wrong preference labels) enters the Evaluator weights directly and is hard to detect after the fact. P6's input noise comes from the Evaluator itself---a miscalibrated Evaluator systematically misleads the Policy.

\subsubsection{Process-Side Dimensions}

The four process-side dimensions form the analytical core of pathway comparison. They characterize the cost, information dynamics, and maintenance properties of the transformation operation itself---what carrier attributes cannot cover at all.

\paragraph{Transformation Semantics.} What kind of re-representation does the transformation operation perform on the experience? This is the most basic characterization of a pathway---different transformation semantics decide how information is organized and which subsequent operations are possible.

\textbf{P1 Narrative Abstraction} performs semantic abstraction: the concrete steps, temporal details, and environment specifics of the raw trajectory are filtered out, and the cross-task transferable rules, patterns, and insights are kept. It does not change the existence level of the information (within the Tokenized level) but changes its organization---from a time-unfolded sequence ($e_1, e_2, \dots, e_T$) to topic- or function-organized entries (a reflection about X, a rule about Y). The cost of this reorganization is that the causal structure of the sequence may be lost---a reflection ``in this scenario one should check B before doing A'' omits the concrete failure process in the raw trajectory by which the rule was discovered, so a future agent can only trust the conclusion and cannot trace back its experiential basis.

\textbf{P2 Schematic Formalization} performs structural formalization: extracting from a natural-language trajectory a syntactic structure a parser or executor can recognize. It does not do semantic abstraction directly---the Mineflayer JavaScript program Voyager generates keeps the concrete actions of the trajectory but weaves them into an executable unit with control flow (loops, conditionals) and parameterization. The gain of formalization is precision: the behavior of a code skill is deterministic (given the input, the output is reproducible), unlike a Narrative reflection, which the model may read differently in different contexts. The cost of formalization is that the ``soft information'' of natural language---implicit assumptions, situational nuance, unstated constraints---may be discarded in the process, because it cannot be captured by a structural grammar.

\textbf{P3 Latent-Space Transformation} performs continuous compression: discrete token sequence $\rightarrow$ continuous vector representation. The distinctive thing about this transformation is that what it keeps is not the semantic content of the experience but the model's internal computation state when processing it---the hidden representations at each layer, attention's organization of the history, the model's understanding after ``digesting'' the experience. These states cannot be fully recovered from the token sequence (otherwise no latent would be needed); they are a lossy complement to the discrete representation. The compressed product has no symbolic readability but can participate in subsequent attention computation directly without re-encoding.

\textbf{P4 Evaluator Internalization} performs the internalization of evaluation capability: interaction experience $\rightarrow$ judgment function. The association between action and outcome in the source experience is learned as a mapping from $(c,a,o)$ to an evaluation signal. The key operation of the transformation is generalization---a trained Evaluator must give an accurate judgment on unseen trajectories too, not merely memorize the $(c,a,o) \rightarrow$ label pairs of the training set. This generalization step is both the source of value (learning evaluation patterns not covered in the training set) and the source of risk (a generalization error causes a systematic misjudgment).

\textbf{P5 Policy Internalization} performs the internalization of decision capability: interaction experience $\rightarrow$ generation function. Its symmetry with P4 is that both are Tokenized $\rightarrow$ Parametric and both rely on gradient-based training. Its asymmetry is that P4 learns ``judgment'' while P5 learns ``execution''---the former's output is an evaluation signal (low-dimensional: a scalar or preference), the latter's an action (high-dimensional: a token sequence or tool call). Learning a high-dimensional output is naturally harder, needs more data, and is less stable.

\textbf{P6 Evaluator-Driven Optimization} performs preference transfer: the Evaluator's judgment capability $\rightarrow$ the Policy's generation preference. The core operation of this transformation is not learning a new capability---the Policy already acquired basic execution capability in the P5 stage---but adjusting the Policy's selection tendency among several feasible actions toward those the Evaluator deems better. The key is signal efficiency: the Evaluator's scalar judgment of a trajectory must be decomposable into a credit assignment for each decision step in the trajectory. The difficulty of this step---the substance of the methodological dispute between RLHF and DPO---is how to recover a step-level learning signal from a trajectory-level evaluation signal.

\textbf{P7 Parametric Externalization} performs knowledge externalization: implicit weight knowledge $\rightarrow$ explicit token sequence. This is the only one of the seven pathways that moves the existence level from inside to outside, the opposite direction of P3/P4/P5. The core mechanism of the transformation is sampling---extracting a concrete behavior instance (trajectory, demonstration, critique) from the Policy's or Evaluator's weights through generation. The fundamental limit of sampling is that what can be extracted from the weights is only a sample of the knowledge they encode, not the whole of it. Different sampling strategies (temperature, prompt, task distribution) produce different externalization results; there is no ``complete externalization.''

\paragraph{Production Mechanism.} What computational infrastructure does the transformation run on? The analytical value of this dimension is that it decides engineering feasibility directly---whether you can afford a pathway depends not on whether it is ``good'' but on whether you have the matching compute resources.

\textbf{LLM inference (prompting)} is the mainstream production mechanism of P1 and P2. The cost model is a token-billed API call: a single transformation's resource need is very low (one trajectory $\rightarrow$ one reflection, usually a few thousand tokens of input plus a few hundred of output). No GPU cluster, no training-data management. The limitation is that transformation quality is bounded by the prompt design and the model's instruction-following ability---you cannot systematically optimize the transformation behavior the way training can.

\textbf{Gradient-based training} is the production mechanism of P3 (Prompt/Module-Based), P4, P5, and P6. The cost model varies across orders of magnitude: SFT a few GPU-hours; a full RLHF pipeline (RM training plus PPO/DPO) tens to hundreds of GPU-hours; the training of a large-scale reasoning model up to thousands of GPU-hours. Beyond raw compute, it also involves training-data management, checkpoint selection, hyperparameter search, and other infrastructure overhead.

\textbf{Gradient-free persistence} is the special mechanism of P3 (Cache-Based): no training, saving the KV cache or hidden state produced by the forward pass directly. The cost is very low---essentially storage cost plus one forward pass. But the price is that it does not cross sessions, does not generalize, and its retrieval efficiency falls as the storage grows.

\textbf{Execution and verification loop} is the additional mechanism of P2 (Programmatic Skill). An LLM-generated code skill must actually execute in the environment and be verified before admission---admitted if it passes, discarded or retried if it fails. The cost of execution verification depends on the environment: a web agent's browser-automation verification is at the second scale, a Minecraft skill's verification may be at the minute scale, and an embodied real-environment verification is slower still. This loop is the production overhead specific to P2 relative to P1---P1 needs no product verification (its verification is the subjective ``reading''), whereas P2 filters out unqualified products through an objective execution signal.

\paragraph{Information Dynamics.} This is the most central and most easily overlooked dimension in pathway comparison. A transformation operation does three things at once---preserve information, lose information, and generate new information---and the weight and nature of these three lines differ sharply across pathways.

\textbf{Information Preservation.} What of the source experience does the transformation keep? P1 keeps the recurring patterns, causal associations, and success/failure conditions of the source trajectory; the incidental details of a single trajectory (the concrete argument value of some tool call, the exact wording of some reasoning) are filtered out, and what is kept is the cross-task transferable semantic content. P2 keeps the process structure and action sequence of the source trajectory: Voyager's JS program keeps the standard operating procedure of crafting, and AWM's~\cite{Wan24} workflow trajectory keeps the phase division and step order of task execution---the procedural knowledge of ``how to do it,'' not the semantic explanation of ``why do it this way.'' P3 keeps the model's internal computation state when processing the source experience---the hidden activations at each layer, attention's key-value pairs---which contain information the token sequence cannot fully encode: the model's implicit tracking of long-range dependencies, its resolution of ambiguity, its dynamic weighting of context. P4 and P5 keep the statistical regularities in the training data---judgment regularities (Evaluator) or decision regularities (Policy); a single experience's contribution is diluted by the accumulated batch gradient and cannot be deconstructed item by item, so what is kept is the overall distribution, not the individual experience. P6 keeps the Evaluator's judgment preference---what behavior is good and what is not---but this passes through two lossy compressions (trajectory to Evaluator in P4, then Evaluator signal to Policy weights in P6), after which the experiential semantics of the original trajectory is no longer identifiable. P7 keeps the behavioral capability encoded in the weights---what it can do---but only as one sample of that capability, not a complete description.

\textbf{Information Loss.} The key distinction is lossy-by-design versus lossy-by-limitation. Lossy-by-design is an intrinsic requirement of the transformation semantics---the loss is deliberate and its category controllable. P1's abstraction necessarily discards detail: distilling from ten failed trajectories of the same task a single rule ``B must be checked before A'' discards each trajectory's specific failure context but keeps the common pattern across instances. P2's formalization necessarily discards the soft information of natural language: a workflow step ``navigate to login page'' does not encode the visual differences of a login page across websites, and these nuances are sacrificed in formalization. Lossy-by-limitation is an engineering limit of the transformation mechanism---the loss is not deliberate and its category uncontrollable. P5 SFT's underfitting may leave some behavior patterns entirely unlearned; RL reward hacking may make the policy learn a shortcut that maximizes the reward signal rather than genuinely improving task performance. P3's KV preservation is bounded by storage capacity---the KV of a long-horizon trajectory must undergo eviction or compression, and which information is discarded is decided by the eviction policy, not by the importance of the experience. The two have sharply different recovery paths on the production side: the former's recovery does not change the transformation method itself (you can keep the original trajectory as a footnote, such as a reflection carrying a reference to its source trajectory), whereas the latter's recovery requires improving the transformation method---better SFT data filtering, reward shaping, a better KV eviction policy---an engineering improvement, not a framework-level design trade-off.

\textbf{Information Gain.} Does the transformation produce new information not explicitly present in the source experience? P1's multi-trajectory induction is the clearest case: comparing ten trajectories can produce statistical knowledge such as ``this kind of failure occurs in 80\% of cases under condition X,'' which comes from no single trajectory but from the cross-trajectory comparison operation itself---a genuine information gain. P2's formalism can also produce a gain: the dependency between steps in a raw trajectory is implicit (the agent always does B before C, but it is never explicitly annotated as a dependency), and formalizing into a workflow DAG encodes it explicitly---not a new fact but the explicitization of an existing one, so downstream can use the dependency mechanically rather than semantically. P3's latent can capture statistical regularities not explicitly encoded by the discrete tokens---for example, some attention head consistently attending to some token position while processing a trajectory, a pattern that comes from the model's way of processing rather than the explicit content. P6 can produce a special gain: the Evaluator's signal may guide the Policy to discover behavior unlearnable from trajectory demonstration alone---if the Evaluator prefers ``a concise and accurate answer,'' the Policy may spontaneously explore in RL a concise expression never present in the demonstrations; this is the power of search plus reward, but also a risk, since the explored ``new behavior'' may be a product of reward hacking. P7's information gain may be negative: a trajectory the Policy externalizes may contain the Policy's hallucination---the sampling noise of the weights externalized as knowledge---which does not constitute ``knowledge'' in the source weights but behaves as experience after externalization.

\textbf{Traceability.} Which parts of the transformation product correspond back to which fragments of the source experience? The analytical value of this dimension is not interpretability itself (the carrier attribute of Section~\ref{sec:carriers}) but that it decides \emph{whether, when experience is audited or corrected externally, the problem lies in the source experience or in the transformation operation}. P1's traceability is the highest: a reflection ``B must be checked before A'' can carry a reference (``see episode \#42 step 7 and episode \#58 step 4''), the structural convenience of a transformation within the Tokenized level. P2's is intermediate: each phase of a workflow can correspond to a time span of the source trajectory, but the correspondence passes through generalization (the phase may span 3 to 7 steps in the source, standardized to 5 in the workflow). P3's drops off a cliff: each dimension of a continuous vector corresponds to no token---probing can find some dimensions' sensitivity to some semantic attributes, but cannot answer ``which token position of which trajectory this KV came from.'' P4 and P5's is near zero: one weight update is the accumulated gradient of all samples in the batch---you can compute the magnitude of a single sample's contribution (through the gradient norm) but cannot locate ``the weight pattern corresponding to this experience'' in the final weights, a structural limit of batch training rather than an engineering-precision problem. P6's is worse than P5's: after the Evaluator $\rightarrow$ Policy transformation, the source of the signal (the original trajectory) is no longer traceable, only ``this Evaluator gave this score in this rollout''---and why the Evaluator gave it returns to P4's untraceability. P7's is special: the product (a tokenized trajectory) is itself readable and traceable, checkable sentence by sentence, but cannot be traced back to the specific components in the weights that produced it---what you see is the Policy's ``behavior,'' not its ``knowledge structure,'' and the provenance chain breaks at P7's sampling step.

\paragraph{Incrementality.} When new experience keeps arriving, can the transformation operation run incrementally? The analytical value of this dimension is that it decides the \emph{sustainability} of an experience-management system---if every new experience requires redoing the transformation from scratch, the system's lifespan is hard-limited by the transformation cost. P1 (Static Store) has the strongest incrementality: new trajectory $\rightarrow$ LLM inference generates a new reflection $\rightarrow$ append to store; existing entries are unaffected, and the overhead is linear in the new experience volume and does not grow with the total. P1 (Dynamic Store) has incrementality with a cascade: a new entry may trigger a merge (combining semantically close ones), prune (evicting low-utility ones), or edit (resolving conflicts) of old entries, but the cascade scope is limited to local entries related to the new one and usually does not trigger a full-store reconstruction, so the overhead is still sub-linear. P2 (Programmatic Skill) is like P1 Static---a new skill is generated, tested, and admitted independently---but P2 (Procedural Workflow) incremental update may trigger a workflow versioning problem: if new experience reveals an erroneous step in an old workflow, one must decide whether to fix the old version or create a new one. P3 (Cache-Based) is naturally incremental: a new session produces new KV cache, the KV of old sessions stays unchanged, and cross-session incrementality shows up as an append to the memory store rather than an update; P3 (Prompt/Module-Based) needs batch retrain---after adding experience the soft prompt or memory module must be retrained, and although it can warm-start from an old checkpoint, the new data changes all the already-learned latent representations, so an incremental update of a single experience is impossible. P4 and P5 have the weakest incrementality: standard preference-model training and SFT/RL need batch retrain---adding training data means either retraining (compute overhead super-linear in the data volume) or fine-tuning (the new data overwrites patterns in the old, with a catastrophic-forgetting risk); online RL (P5 and P6) can update batch by batch, but each update uses only the current batch's information and the effect of old-batch experience decays over time---not genuine ``incrementality'' but ``forgetting old experience, relying only on recent experience.'' P7's incrementality depends on the source (Policy or Evaluator): if the Policy is continually updated through online RL, each P7 externalization reflects the latest Policy state; if the Policy is batch-retrained, P7 runs in batch mode too. Incrementality connects directly to the temporal-evolution analysis of Section~\ref{sec:pathway-selection}: if a system is expected to run long-term and experience keeps growing, choosing a non-incremental pathway (P4/P5) means either accepting the compute overhead of periodic retraining or accepting the decay of experience utilization over time---a hard constraint on scenario recommendation.

\paragraph{Output Verifiability.} After the transformation completes, can one independently verify whether the transformation operation itself succeeded? Note this is not the interpretability of the target carrier (a consumption-side issue, analyzed in Section~\ref{sec:utilization-analysis}) but \emph{whether the product's quality can be judged without relying on downstream task performance}. P1's Output Verifiability is direct but subjective: a generated reflection can be read by a human or a strong LLM to judge whether the content is reasonable, consistent with the source trajectory, and free of obvious errors---a qualitative check, not a quantitative test, with no objective failure criterion (a reflection can be ``roughly correct but misleading'' and still be judged usable). P2 (Programmatic Skill)'s is objective and automated: a generated code skill is executed in the environment, and pass/fail is a binary signal whose pass rate is countable; but test completeness cannot be guaranteed (a skill passing all existing tests may fail on an uncovered corner case)---a coverage problem, not an objectivity problem. P2 (Workflow/Graph)'s is intermediate: structural correctness can be checked through static analysis (all step dependencies satisfied, no dead loops, input/output types matched), but task completeness (whether all necessary processing branches are covered) cannot be verified automatically and needs human judgment or an end-to-end test. P3's is the weakest: a latent representation's quality has no independent evaluation criterion---you can only infer it indirectly through ablation (the performance change after removing this latent bank), and that inference confounds the quality of the representation with the way it is used (the same high-quality KV cache, selected wrongly by an inappropriate retrieval strategy, also performs poorly in ablation). P4's is weak and indirect: an Evaluator's quality is usually evaluated through two levels of proxy---agreement with human labels (themselves noisy, with a distribution that may differ from the Evaluator's training distribution) and utility in downstream P6 (where Policy training is affected by many factors and the Evaluator quality's contribution cannot be isolated). P5's can only be evaluated indirectly through rollout (whether the new policy outperforms the old on hold-out tasks), but this confounds several factors---the quality of the new trajectories, the training hyperparameters, the old policy's performance ceiling---and a policy that ``has not regressed'' may have been updated dominated by noise rather than by an effective fixing of experience. P6 faces a circular-verification risk: using the Evaluator's score to train the Policy and then the Policy's performance to verify the Evaluator's quality, the two referencing each other with no independent external ground truth, which compounds the already weak verification of P4 and P5. This dimension connects directly to the structural regularity verifiability $\rightarrow$ correctability $\rightarrow$ maintainability: if you cannot independently verify whether a transformation succeeded, you cannot judge when correction is needed, and so cannot maintain experience quality over long-term operation. In the selection logic of Section~\ref{sec:pathway-selection}, a path low in Output Verifiability is at a systematic disadvantage in scenarios that need long-term maintenance.

\subsubsection{Output-Side Dimension: Composability}

Can a pathway's product serve directly as the input of another pathway? This dimension characterizes a pathway's out-degree in the directed graph of the seven pathways---it decides the experience's ``fluidity'' in the carrier system.

The difference in Composability is structural rather than empirical---it is entailed by the functional-role principle (Section~\ref{sec:functional-role-principle}). Policy produces actions, actions produce trajectories, and a trajectory can be consumed by any pathway sourced in Tokenized---the Policy's Composability comes from its functional definition, not from accidental community exploration. Evaluator produces evaluation signals whose semantics is ``a judgment of behavior,'' and the only consumer is the producer of behavior---the Evaluator's Composability of $1$ (only P6) is likewise a structural corollary of the functional definition.

P1 (Narrative) is the node of highest out-degree. A Narrative artifact (reflection, rule, insight) can:
\begin{itemize}
  \item serve as input to P2 (extracting a formalizable operation process from a reflection);
  \item serve as input to P3 (compressing a reflection into a latent for fast context injection);
  \item serve as input to P4 (training an Evaluator from the comparison of a reflection and the trajectory pair it corrects);
  \item serve as input to P5 (a reflection as a supplementary or substitute supervision signal for a trajectory);
  \item be consumed directly (injected into the context as a retrieval result).
\end{itemize}

P2 (Schematic) has medium out-degree. The main downstreams of a Schematic artifact are:
\begin{itemize}
  \item input to P5 (a skill as an extension of the action space, or a skill execution trace as training data);
  \item direct execution consumption.
\end{itemize}

P3 (Latent) currently has almost no literature as the source end of a new pathway. This reflects a deep asymmetry: Tokenized $\rightarrow$ Latent has a clear motivation (preserving the computation state, lowering the context overhead), but Latent $\rightarrow$ Tokenized or Latent $\rightarrow$ Parametric lacks a symmetric one---if you want to turn a latent back into tokenized, you may as well keep the original tokenized experience; if you want to fix a latent into the weights, that is a variant of P5 rather than an independent new pathway. The Latent's role in the pathway graph is a ``way station''---experience enters and usually stops in latent form for consumption, rather than continuing to transform from there.

P4 (Evaluator)'s out-degree is $1$: only into P6, as argued in the functional-role principle.

P5 (Policy)'s out-degree is $2$: into P7 (explicit externalization), and through rollout implicitly into any pathway sourced in Tokenized. The Policy's Composability makes it the hub of composite pathways---the self-generation loop (Policy $\rightarrow$ rollout $\rightarrow$ Tokenized $\rightarrow$ P5) and Evaluator--Policy Co-Evolution (Policy $\rightarrow$ rollout $\rightarrow$ P4 $\rightarrow$ P6 $\rightarrow$ Policy) both rely on this one structural property.

P6 (Evaluator-Driven Optimization)'s product is a Policy, so its out-degree equals P5's.

P7 (Parametric Externalization)'s product is a Tokenized trajectory or demonstration, so its out-degree equals P1's---in principle it can enter any pathway sourced in Tokenized. P7 is the loop-closing node of the pathway graph: \textbf{Policy $\rightarrow$ P7 $\rightarrow$ Tokenized $\rightarrow$ P5 $\rightarrow$ Policy forms the self-generation loop}, the structural basis of the Generative Experience Curation of Section~\ref{sec:experience-curation}.

Table~\ref{tab:pathway-comparison} summarizes the comparison of the seven pathways across seven dimensions.

\begin{table*}[t]
\centering
\caption{Comparison of the seven pathways across seven dimensions. Rows are pathways, not methods; cells state structural constraints, not literature counts.}
\label{tab:pathway-comparison}
\scriptsize
\setlength{\tabcolsep}{2.5pt}
\renewcommand{\arraystretch}{1.15}
\begin{tabular}{@{}p{1.5cm}p{0.55cm}p{2.6cm}p{2.1cm}p{2.2cm}p{1.9cm}p{2.2cm}p{2.1cm}@{}}
\toprule
Pathway & Source & Data Prerequisite & Production Mechanism & Info.\ Loss Type & Traceability & Incrementality & Output Verifiability \\
\midrule
P1 Narrative Abstraction & $\mathcal{C}^T_N$ & Low volume; no labels needed; self-gen or human & LLM inference & Lossy-by-design (abstraction) & High (artifact $\rightarrow$ source episode) & Incremental (append) / with merge (dynamic) & Direct but subjective (reading) \\
\addlinespace
P2 Schematic Formalization & $\mathcal{C}^T_N$ & Low--medium volume; no labels needed; self-gen or human & LLM inference + execution verification & Lossy-by-design (formalization) & Medium (structure $\rightarrow$ trajectory segment) & Incremental (append skill) & Objective but coverage-limited (execution test) \\
\addlinespace
P3 Latent-Space & $\mathcal{C}^T$ & Session: minimal; trained: large volume; labels only for Module-Based & Forward pass + KV persistence / gradient training & Lossy-by-design (compression) & Low (vector $\rightarrow$ batch) & Session: per-session; trained: batch retrain & Weak (ablation only) \\
\addlinespace
P4 Evaluator Internalization & $\mathcal{C}^T$ & Large volume; outcome or process labels; human or teacher & Gradient-based training & Mixed (compression by-design + training loss by-limitation) & Very low (weight update $\rightarrow$ batch) & Batch / full retrain & Weak, indirect (human agreement or P6 utility) \\
\addlinespace
P5 Policy Internalization & $\mathcal{C}^T$ & Large volume; outcome-filtered labels or none (RL); self-gen or human & Gradient-based training & Mixed (compression by-design + underfitting / reward hacking by-limitation) & Very low (weight update $\rightarrow$ batch) & Batch retrain / online iterative & Behavioral only (rollout comparison) \\
\addlinespace
P6 Evaluator-Driven Opt. & $\mathcal{C}^P_\phi$ & Medium--large volume (Evaluator-dependent); Evaluator signals as implicit labels & Gradient-based training (RL / DPO) & Double compression (P4 loss + P6 loss) & Near-zero (two steps from source) & Online iterative & Circular risk (Evaluator $\leftrightarrow$ Policy) \\
\addlinespace
P7 Parametric Externalization & $\mathcal{C}^P$ & No training; sampling quality depends on prompt design & Sampling from trained model & Lossy-by-design (sampling) + hallucination risk & Medium (artifact readable, not traceable to weights) & Follows source-model update & Artifact readable, completeness unverifiable \\
\bottomrule
\end{tabular}
\end{table*}

This table is not a summary of Sections~\ref{sec:tok2tok}--\ref{sec:param-source}---its rows are pathways, not methods, and the cells contain structural constraints, not literature counts. Take P1's Incrementality: the table says ``Incremental (append) / with merge (dynamic),'' the two values corresponding to the structural difference between the Static Store and Dynamic Store classes in Section~\ref{sec:narrative-abstraction}, not the names of specific works. A reader who wants to know which works are Static and which Dynamic should return to the table and prose of Section~\ref{sec:narrative-abstraction}.

Table~\ref{tab:pathway-comparison} is the output of Section~\ref{sec:pathway-tradeoff} and the data basis of the production-path-selection phase of Section~\ref{sec:pathway-selection}. The three cross-section dimensions---input side, process side, output side---each have a role: the input side decides what quality of raw material you can supply this path; the process side decides whether you can afford the computation and labor of the transformation itself; the output side decides the product's fluidity in the subsequent experience-management system.

\subsection{Carrier Utilization Analysis}
\label{sec:utilization-analysis}

% Content forthcoming.

\subsection{Utilization-Driven Pathway Selection}
\label{sec:pathway-selection}

% Content forthcoming.
